% =============================================================================
% Section VI — Discussion
% =============================================================================
\section{Discussion}
\label{sec:discussion}

The current model demonstrates a meaningful ranking signal (AUC-ROC 0.6364,
AUC-PR 0.5008) but an imbalanced operating point. With $\tau=0.0913$, recall is
perfect on the test split while precision and accuracy are low. In practical
terms, the detector is conservative with respect to missed leaks and aggressive
in generating alarms.

\subsection{What Worked}
\begin{itemize}[leftmargin=*]
	\item Multivariate sequence modeling (CNN + BiLSTM + attention) captures leak-related
		  temporal structure beyond static thresholding.
	\item Feature engineering supports class ranking quality, as reflected in PR/ROC curves.
	\item The evaluation pipeline provides transparent threshold diagnostics for deployment.
\end{itemize}

\subsection{Main Limitation}
The deployed threshold is too low for balanced operation on this test set.
As shown by threshold sensitivity plots, moderate threshold increases can
substantially reduce false alarms, but this may decrease recall.

\subsection{Deployment Implication}
Threshold should be selected from an explicit operational objective, for example:
\begin{itemize}[leftmargin=*]
	\item Recall-constrained optimization (maximize precision with recall $\geq r_{min}$), or
	\item Cost-based optimization using field-inspection and missed-leak penalties.
\end{itemize}

\subsection{Future Work}
\begin{itemize}[leftmargin=*]
	\item Re-train and export model without Lambda deserialization fragility by replacing
		  Lambda reduction with a native pooling layer.
	\item Run cross-validation across temporal folds to improve robustness of reported metrics.
	\item Extend evaluation with calibration curves and uncertainty-aware alert policies.
\end{itemize}

